% TOP_PROBLEM: {"id": 2764, "title": "Kirby Problem 2.16", "category_id": 7, "category": {"id": 7, "name": "topology", "display_name": "Topology", "description": "Properties preserved under continuous deformations.", "slug": "topology", "order_index": 7, "created_at": "2026-07-31T15:26:25.671Z"}, "status": "open", "problem_number": "KP-2.16", "set_id": 11}
% FIRST_POSTED: 2026-10-01
% ENTRY_KIND: statement_only
% UnsolvedMath ID 2764; source catalogue code KP-2.16
% !TeX program = lualatex
% Definition and sources only; this file is not an LLM solution attempt.
\documentclass[11pt,a4paper]{article}
\usepackage[margin=25mm]{geometry}
\usepackage{fontspec}
\setmainfont{lmroman10-regular.otf}[BoldFont=lmroman10-bold.otf,
 ItalicFont=lmroman10-italic.otf,BoldItalicFont=lmroman10-bolditalic.otf]
\usepackage{amsmath,amssymb,mathtools}
\usepackage{unicode-math}
\setmathfont{latinmodern-math.otf}
\AtBeginDocument{\let\setminus\smallsetminus}
\usepackage{xurl,hyperref}
\hypersetup{colorlinks=true,urlcolor=blue}
\setcounter{secnumdepth}{0}
\setlength{\parindent}{0pt}
\setlength{\parskip}{5pt}
\setlength{\emergencystretch}{3em}
\begin{document}
\section*{Kirby Problem 2.16}\label{2764}
{\small UnsolvedMath ID 2764 · KP-2.16 · Topology · Kirby's Problems in Low-Dimensional Topology}
\subsection{Definitions and mathematical statement}
For a complete finite-area hyperbolic surface $X$ without boundary, each
essential nonperipheral
free homotopy class of simple closed curves has a unique closed geodesic
representative. Its unmarked simple length spectrum is the multiset
\[
\mathcal L_{\rm simp}(X)=
\{\!\{\ell_X(\alpha):\alpha\text{ an unoriented simple closed
geodesic class}\}\!\}.
\]
Count one entry for each unoriented class; distinct classes of equal length
produce multiplicities. Iterated traversals of a primitive geodesic are not
additional simple curves. No labels identifying the classes on a reference
surface are retained.

If two hyperbolic surfaces $X,Y$ have equal unmarked simple length spectra,
must they be isometric? In particular, ask this for closed orientable surfaces.
An isometry may preserve or reverse orientation, since lengths alone contain no
orientation choice.
Equality means matching the multiplicity of every real length, rather than
merely equality of the sets of values. The challenge is to recover the metric
from the unlabeled data; the hypothesis does not supply a homeomorphism that
matches individual curve classes. Such a supplied marking would give a
different, stronger hypothesis.
\subsection{Short English statement}
Does the multiset of lengths of simple closed geodesics determine a hyperbolic surface up to isometry?
\subsection{Sources}
\begin{itemize}
\item[S1] ULAM.AI and the UnsolvedMath Contributors, supplied problems.json, record 2764 (KP-2.16), Kirby's Problems in Low-Dimensional Topology. Catalogue identity and source transcription; CC BY 4.0.
\url{https://huggingface.co/datasets/ulamai/UnsolvedMath}.
\item[S2] K3: A New Problem List in Low-Dimensional Topology, Problem 2.16. Expanded formulation of the supplied catalogue statement.
\url{https://math.berkeley.edu/sites/default/files/surv-295-ruberman-watermarked-author-pdf.pdf}.
\end{itemize}
\end{document}
